% Einheit 2 von 4 – Hauptsatz (bank/stammfunktion-und-hauptsatz/e2.jsonl, zusammenbau v0.1)
%% TODO Zeitmarke Einheit 2: Marken-Zeile nicht lesbar
%% TODO Zweigzeile Teil 1: was hier gelernt wird (Satz in Schülersprache aus dem Einheitstitel) – Einheit 2
\einheitenkopf[e2]{Einheit 2 von 4 · Hauptsatz}
\zweigzeile{Abitur GK}
\setcounter{aufgabe}{9}

%% TODO Titel: Ich-kann-Satz fehlt in der Bank (Platzhalter: Kettenname) – Nr. 10
%% TODO Anweisung: ein Satz über der ersten Teilaufgabe fehlt in der Bank – Nr. 10
\begin{aufgabe}{Hauptsatz}
%% TODO \rechenplatz{n}: Schrittzahl des Grundfalls fehlt in der Bank (nur bei fester Schreibform, 2.3 b)
\begin{teile}
\teil Berechne das Integral von $0$ bis $2$ über $f(x) = 3x^2 - 2x$. Wert: \leerfeld
\teil $F(x) = 3x - 4e^{-\frac{x}{2}}$ ist eine Stammfunktion von $f$. Berechne das Integral von $0$ bis $2$ über $f(x)$ auf zwei Stellen. \hfill (Abitur 2026 LK) \\ Wert: \leerfeld
\teil Berechne das Integral von $0$ bis $2\pi$ über $f(x) = \mathrm{cos}(x) + 2$. \hfill (Abitur 2021 GK) \\ Wert: \leerfeld
\teil $g(x) = (x - 1)^2 \cdot e^{x}$. Berechne das Integral von $0$ bis $1$ über $g'(x)$. \hfill (Abitur 2025 LK) \\ Wert: \leerfeld
\teil $F(x) = (2 - x) \cdot e^{x}$ ist eine Stammfunktion von $f(x) = (1 - x) \cdot e^{x}$. Ein Näherungswert für das Integral von $-2$ bis $1$ über $f(x)$ ist $2{,}1$. Berechne den exakten Wert und die prozentuale Abweichung des Näherungswerts. \hfill (Abitur 2025 GK) \\ exakt: \leerfeld , Abweichung: \leerfeld \%
\end{teile}
\end{aufgabe}

%% TODO Titel: Ich-kann-Satz fehlt in der Bank (Platzhalter: Kettenname) – Nr. 11
%% TODO Anweisung: ein Satz über der ersten Teilaufgabe fehlt in der Bank – Nr. 11
\begin{aufgabe}{Kurvenlänge über eine vorgegebene Integralformel berechnen}
\begin{teile}
\teil Für die Länge $s$ des Graphen von $f$ zwischen $a$ und $b$ gilt: $s$ ist das Integral von $a$ bis $b$ über $\sqrt{1 + (f'(x))^2}$. Berechne die Länge des Graphen von $f(x) = \frac{2}{3}\sqrt{x^3}$ zwischen $x = 0$ und $x = 3$. Eine Stammfunktion von $\sqrt{1 + x}$ ist $\frac{2}{3}\sqrt{(1 + x)^3}$. $s = $ \leerfeld
\end{teile}
\end{aufgabe}

%% TODO Titel: Ich-kann-Satz fehlt in der Bank (Platzhalter: Kettenname) – Nr. 12
%% TODO Anweisung: ein Satz über der ersten Teilaufgabe fehlt in der Bank – Nr. 12
\begin{aufgabe}{Hauptsatz – Fehler finden, Begründen}
\begin{teile}
\teil Ich finde den Fehler: Mia berechnet das Integral von $1$ bis $3$ über $2x$ so: $F(1) - F(3) = 1 - 9 = -8$.
\teil Begründe, warum man für das Integral von $1$ bis $4$ über $f$ jede Stammfunktion von $f$ nehmen darf.
\end{teile}
\end{aufgabe}
